\section{A finite-size criterion for the cube condition}\label{app:finite}

This appendix proves the finite-size criterion used in
Sec.~\ref{sec:massgen}: a gap for the dynamics that resamples the blocks of
one patch of $m^4$ blocks, uniform in the configuration outside the patch,
gives the cube condition on every large torus. Blocks that share no
plaquette give commuting projections, and this replaces the disjointness
used in the finite-size criteria of Knabe~\cite{Kna88} and of Gosset and
Mozgunov~\cite{GM16} for quantum spin systems.\footnote{In those criteria
the local terms of a Hamiltonian act on overlapping sets of sites, and the
gap of the whole system is bounded below by the gap of a finite patch minus
a correction that decays with the size of the patch. Here the local terms
are the complements $I-\mathcal P_j$ of conditional expectations, and two
of them commute as soon as the blocks are at distance at least two.}

\enlargethispage{5pt}
\emph{Setting.} The notation is that of Sec.~\ref{sec:cube} and of the
setting of Sec.~\ref{sec:massgen}. Fix integers $d\ge2$ and $m\ge4$ with
$d\mid L$, put $K=L/d$, $n=2d-1$, and fix a shift $a\in[0,d)^4$. For $j$ in
the coarse torus $T_K=(\Z/K\Z)^4$ let $r_j=a+dj$, let $B_j=V_{r_j}$ be the
set of links with base point in $r_j+[0,n)^4$, and write $\mathcal P_j=
\mathcal P_{r_j}$. Displacements $v=j'-j$ in $T_K$ are represented in
$(-K/2,K/2]^4$. Two sets $A,A'$ of links are \emph{non-interacting} if no
plaquette contains a link of each; for a set $A$ of links, $P_A$ is the
conditional expectation under $\mu$ given the links outside $A$.

\begin{lemma}\label{lem:fcommute}
\begin{enumerate}[label=\textup{(\roman*)},leftmargin=*]
\item For non-interacting sets $A,A'$ of links, $P_AP_{A'}=P_{A'}P_A=
P_{A\cup A'}$, and $(I-P_A)(I-P_{A'})$ is an orthogonal projection.
\item If $K\ge4$ and $|j'-j|_\infty\ge2$, then $B_j$ and $B_{j'}$ are
non-interacting.
\item The blocks $B_j$, $j\in T_K$, cover $E$.
\item If $f\in L^2(\mu)$ and $\mathcal P_jf=f$ for every $j\in T_K$, then
$f$ is almost surely constant.
\end{enumerate}
\end{lemma}

\begin{proof}
(i) Every link lies on a plaquette, so $A\cap A'=\emptyset$. Given the
links outside $A\cup A'$, the law of $(U_A,U_{A'})$ has density
proportional to $\prod_{p\cap(A\cup A')\ne\emptyset}\varphi(U_p)$, and each
factor involves the links of exactly one of $A$, $A'$. The density is
therefore a product of a function of $U_A$ and a function of $U_{A'}$,
with coefficients depending on the links outside $A\cup A'$. Hence $U_A$
and $U_{A'}$ are conditionally independent, the conditional law
$\gamma_A$ of $U_A$ given the links outside $A$ does not depend on
$U_{A'}$, and likewise for $A'$. So $P_AP_{A'}f$ and $P_{A'}P_Af$ are both
the integral of $f$ against $\gamma_A\otimes\gamma_{A'}$, which is
$P_{A\cup A'}f$, and $(I-P_A)(I-P_{A'})=I-P_A-P_{A'}+P_{A\cup A'}$ is
self-adjoint and idempotent.

(ii) Let $|v_i|\ge2$ for $v=j'-j$. In the coordinate $i$ the base points of
the links of $B_j$ and of $B_{j'}$ lie in $dj_i+a_i+[0,2d-2]$ and in
$dj_i+a_i+dv_i+[0,2d-2]$ modulo $L$. If $v_i\ge2$, their distance in
$\Z/L\Z$ is at least $\min\{(v_i-2)d+2,(K-v_i-2)d+2\}\ge2$, since $v_i\le
K/2$ and $K\ge4$; the case $v_i\le-2$ is symmetric. A plaquette meeting both
blocks would contain two links whose base points differ by at least two in
the coordinate $i$, while the base points of the links of a plaquette
differ by at most one in each coordinate.

(iii) The sets $r_j+[0,d)^4$ partition $\Lambda_L$, and $[0,d)\subset
[0,n)$.

(iv) The density of $\mu$ with respect to the Haar measure $h$ of $G^E$ is
bounded above and below by positive constants, so $\mu$ and $h$ have the
same null sets and $L^2(\mu)=L^2(h)$. If $\mathcal P_jf=f$, then $f$ is
almost surely a function of $U_{E\setminus B_j}$, and it is fixed by the
conditional expectation $Q_j$ of $h$ given $U_{E\setminus B_j}$, which
integrates out $U_{B_j}$. By Fubini's theorem the $Q_j$ commute and their
product integrates out the links of $\bigcup_jB_j=E$, so $f=\int f\,dh$
almost surely.
\end{proof}

In four dimensions the displacements $v$ with $|v|_\infty=1$ are $80$ in
number, and those with $|v|_\infty\ge2$ give commuting projections. For
$s\in T_K$ let $S_s=s+[0,m-1]^4$ be the \emph{patch} at $s$ and $S^\circ_s=
s+[0,m-2]^4$ the set of its $(m-1)^4$ block indices, and put
\[
H=\sum_{j\in T_K}(I-\mathcal P_j),\qquad H_s=\sum_{j\in S^\circ_s}(I-
\mathcal P_j).
\]
These are bounded nonnegative operators on $L^2(\mu)$ and preserve the
subspace $\mathcal A$ of gauge-invariant functions.

\begin{lemma}\label{lem:knabe}
Let $m\ge4$, $p=m-1$, $K\ge2m$, let $\mathcal K$ be $L^2(\mu)$ or
$\mathcal A$, and suppose that $H_s^2\ge\gamma H_s$ on $\mathcal K$ for
every $s\in T_K$, with $\gamma>0$. Then $H^2\ge\varepsilon H$ on
$\mathcal K$ with
\[
\varepsilon=1-\frac{(1-\gamma)p^4+72p^3-168p^2+96p-16}{p^3(p-2)}\ge
\frac{\gamma p-74}{p-2}.
\]
In particular $\varepsilon\ge\gamma/2$ if $m-1\ge148/\gamma$.
\end{lemma}

\begin{proof}
All operators below preserve $\mathcal K$, and a restriction of an
orthogonal projection to an invariant closed subspace is an orthogonal
projection; so it suffices to argue on $L^2(\mu)$ with inequalities read on
$\mathcal K$. Write $\mathbf P_j=I-\mathcal P_j$ and $A_v=\sum_j\mathbf P_j
\mathbf P_{j+v}$. Since $K\ge2m$, two indices $j$ and $j+v$ lie in a common
$S^\circ_s$ only if $|v_i|\le m-2$ for all $i$, and then for exactly
$c_v=\prod_{i=1}^4(p-|v_i|)$ points $s$. Let $O$ be the set of the $80$
displacements with $|v|_\infty=1$, $D$ the set of those with
$|v|_\infty\ge2$, and $D'$ the set of $v\in D$ with all $|v_i|\le m-2$. For
$v\in D$, $\mathbf P_j\mathbf P_{j+v}$ is an orthogonal projection by
Lemma~\ref{lem:fcommute}(i),(ii), so $A_v\ge0$. Summing the hypothesis over
$s$,
\[
\sum_sH_s^2=c_0H+\sum_{v\in O}c_vA_v+\sum_{v\in D'}c_vA_v\ge\gamma\sum_sH_s
=\gamma c_0H .
\]
Put $c_*=p^3(p-2)$; then $c_v\le c_*$ for $v\in D'$, since one factor of
$c_v$ is at most $p-2$ and the others are at most $p$. For orthogonal
projections $\mathbf P,\mathbf P'$, $(\mathbf P\pm\mathbf P')^2\ge0$ gives
$-(\mathbf P+\mathbf P')\le\mathbf P\mathbf P'+\mathbf P'\mathbf P\le
\mathbf P+\mathbf P'$, so $-2H\le A_v+A_{-v}\le2H$. As $O$ is symmetric and
$c_{-v}=c_v$,
\[
\sum_{v\in O}(c_*-c_v)A_v\ge-\Sigma H,\qquad\Sigma=\sum_{v\in O}|c_*-c_v|,
\]
and therefore
\[
c_*\sum_{v\in O}A_v=\sum_{v\in O}c_vA_v+\sum_{v\in O}(c_*-c_v)A_v\ge
-\bigl((1-\gamma)c_0+\Sigma\bigr)H-\sum_{v\in D'}c_vA_v .
\]
Since $H^2=\sum_vA_v=H+\sum_{v\in O}A_v+\sum_{v\in D}A_v$,
\[
H^2\ge\Bigl(1-\frac{(1-\gamma)c_0+\Sigma}{c_*}\Bigr)H+\sum_{v\in D'}\Bigl(1-
\frac{c_v}{c_*}\Bigr)A_v+\sum_{v\in D\setminus D'}A_v ,
\]
and the last two sums are nonnegative. It remains to compute $\Sigma$. A
displacement $v\in O$ with $k$ nonzero coordinates has $c_v=p^{4-k}(p-1)^k$,
and there are $8$, $24$, $32$ and $16$ of them for $k=1,2,3,4$. The
differences $c_v-c_*$ are
\[
p^3,\qquad p^2,\qquad-p(p^2-3p+1),\qquad-2p(p-1)(p-2)+1,
\]
with the signs shown since $p\ge3$. Hence
$\Sigma=8p^3+24p^2+32p(p^2-3p+1)+16(2p(p-1)(p-2)-1)=72p^3-168p^2+96p-16$,
and with $c_0=p^4$ this is the formula for $\varepsilon$. As $-168p^2+96p
-16<0$, $\varepsilon\ge1-((1-\gamma)p+72)/(p-2)=(\gamma p-74)/(p-2)$, and if
$\gamma p\ge148$ this is at least $(\gamma p-74)/p\ge\gamma/2$.
\end{proof}

For a patch $S_s$ let $\Pi_s=\bigcup_{j\in S^\circ_s}B_j$, the set of links
with base point in $r_s+[0,md-1)^4$. For $L\ge md+2$ the patch and the
plaquettes meeting it embed in $\Z^4$, and the conditional law of
$U_{\Pi_s}$ under $\mu$ given the links outside $\Pi_s$ is the measure
$\nu^\omega_{Q}$ of Sec.~\ref{sec:massgen} for the box $Q=r_s+[0,md-1)^4$,
with $\omega$ the configuration outside. For $j\in S^\circ_s$ write
$\mathcal P^\omega_j$ for the conditional expectation under $\nu^\omega_Q$
given the links of $\Pi_s\setminus B_j$.

\begin{definition}\label{def:patchgap}
The \emph{gauge-invariant patch gap} $\gamma_{d,m}$ is the largest
$\gamma\ge0$ such that
\[
\gamma\Var_{\nu^\omega_Q}(F)\le\sum_{j\in S^\circ_s}\nu^\omega_Q\bigl((F-
\mathcal P^\omega_jF)^2\bigr)
\]
for every patch box $Q=r+[0,md-1)^4$ of $\Z^4$, every exterior
configuration $\omega$ and every $F\in\mathcal A^\omega_Q$.
\end{definition}

By translation invariance $\gamma_{d,m}$ depends only on $G$, $\varphi$,
$d$ and $m$. It is the gap, on gauge-invariant functions, of the dynamics
that resamples the $(m-1)^4$ blocks of one patch, uniform in the
configuration outside the patch.

\begin{theorem}[finite-size criterion]\label{thm:finite}
Let $d\ge2$, $m\ge4$, $d\mid L$, $L$ even, $L\ge2md$, and suppose that
$\gamma_{d,m}\ge\gamma>0$ and that the number $\varepsilon$ of
Lemma~\ref{lem:knabe} is positive. Then
\[
\frac{\varepsilon}{16}\Var_\mu(f)\le\frac1{n^4}\sum_{x\in\Lambda_L}\mu\bigl(
(f-\mathcal P_xf)^2\bigr)\qquad\text{for every real }f\in\mathcal A,
\]
that is, $\BP(2d-1,\varepsilon/16)$. If $m-1\ge148/\gamma$, then
$\varepsilon/16\ge\gamma/32$.
\end{theorem}

\begin{proof}
Fix $a$. Since $L\ge2md$, $K\ge2m\ge8$.

\emph{The patch inequality on the torus.} Let $F\in\mathcal A$; replacing $F$
by $\int F(gU)\,dg$, the integral over all gauge transformations of
$\Lambda_L$ with Haar measure, we may take $F$ invariant at every point.
Write $U=(\sigma,\omega)$ with $\sigma$ the links of $\Pi_s$ and $\omega$ the
others. Under $\mu$ the conditional law of $\sigma$ given $\omega$ is
$\nu^\omega_Q$, and for $j\in S^\circ_s$ the conditional law of $U_{B_j}$
given the other links is the same under $\mu$ and under $\nu^\omega_Q$, so
$(\mathcal P_jF)(\cdot,\omega)=\mathcal P^\omega_j(F(\cdot,\omega))$ for almost
every $\omega$ and $H_s$ acts on each fibre as $H^\omega_s=\sum_{j\in
S^\circ_s}(I-\mathcal P^\omega_j)$. A gauge transformation of $\Lambda_L$
that is trivial at the vertices not interior to $Q$ fixes $\omega$, so each
fibre $F(\cdot,\omega)$ lies in $\mathcal A^\omega_Q$. On $\mathcal
A^\omega_Q$ the operator $H^\omega_s$ is self-adjoint and nonnegative,
preserves the constants and their orthogonal complement
(Lemma~\ref{lem:boxinv}(iii)), and by Definition~\ref{def:patchgap} is at
least $\gamma$ on that complement. So its spectrum there lies in
$\{0\}\cup[\gamma,\infty)$ and $(H^\omega_s)^2\ge\gamma H^\omega_s$.
Integrating over $\omega$, $\|H_sF\|^2\ge\gamma\ip F{H_sF}$ for $F\in
\mathcal A$, that is, $H_s^2\ge\gamma H_s$ on $\mathcal A$.

\emph{Conclusion.} By Lemma~\ref{lem:knabe} with $\mathcal K=\mathcal A$,
$H^2\ge\varepsilon H$ on $\mathcal A$, so the spectrum of $H$ on $\mathcal
A$ lies in $\{0\}\cup[\varepsilon,\infty)$. If $F\in\mathcal A$ and $HF=0$,
then $\sum_j\|F-\mathcal P_jF\|^2=\ip F{HF}=0$, so $F$ is constant by
Lemma~\ref{lem:fcommute}(iv). Hence for real $f\in\mathcal A$ the function
$g=f-\mu(f)\in\mathcal A$ is orthogonal to the kernel of $H$ on $\mathcal
A$, and
\[
\varepsilon\Var_\mu(f)\le\ip g{Hg}=\sum_{j\in T_K}\mu\bigl((f-\mathcal
P_{r_j}f)^2\bigr).
\]
As $a$ runs over $[0,d)^4$ and $j$ over $T_K$, $r_j=a+dj$ runs over
$\Lambda_L$ once. Averaging over $a$ gives $\varepsilon d^4\Var_\mu(f)\le
\sum_x\mu((f-\mathcal P_xf)^2)$, and $d^4\ge n^4/16$. Since $L\ge8d\ge2n+4$
this is $\BP(2d-1,\varepsilon/16)$, and the last claim is
Lemma~\ref{lem:knabe}.
\end{proof}

The passage from the patch to the torus costs a factor: at $\varphi\equiv1$
the links are independent, $\gamma_{d,m}\ge1$ by the Efron--Stein
inequality~\cite{ES81}, and the theorem gives $\BP(2d-1,\varepsilon/16)$
with $\varepsilon<1$, against $\BP(n,1)$ from Proposition~\ref{prop:bpsup}.

The hypothesis of Theorem~\ref{thm:finite} is a statement about every
boundary condition of the patch. It is in fact equivalent to the local
inequality $H_s^2\ge\gamma H_s$ on the torus, so that it is not a stronger
hypothesis than the one the proof uses, and it need only be imposed on the
functions that are even under the centre. Let $\mathcal G_{\rm ext}$ be the
group of gauge transformations of $\Lambda_L$ that are trivial at the
vertices interior to $Q$; it fixes the links of $\Pi_s$ with both end
points interior, acts on the other links of $\Pi_s$ and on $\omega$, and
commutes with $\mathcal G_Q$, and every gauge transformation of $\Lambda_L$
is a product of one element of each. For $x\in G$ let $c_x\in\mathcal G_{\rm
ext}$ be equal to $x$ at every non-interior vertex; it acts on $\omega$ by
conjugation, so $c_z\omega=\omega$ for $z$ in the centre $Z$ of $G$, while
$c_z$ multiplies the links of $\Pi_s$ with one interior end point by $z^{\pm
1}$. Write $\operatorname{Stab}(\omega)=\{g\in\mathcal G_{\rm ext}:g\omega=
\omega\}\ni c_z$, $H^\omega=\sum_{j\in S^\circ_s}(I-\mathcal P^\omega_j)$ on
$L^2(\nu^\omega_Q)$ and
\[
\Phi_\omega(F)=\|H^\omega F\|^2-\gamma\ip F{H^\omega F}\qquad(F\in L^2(\nu^
\omega_Q)).
\]
A function $F\in\mathcal A^\omega_Q$ is \emph{$Z$-even} if $F\circ c_z=F$ for
all $z\in Z$.

\begin{proposition}[the patch inequality fibre by fibre]\label{prop:fibre}
Let $d\ge2$, $m\ge4$, $d\mid L$, $L$ even, $L\ge2md$, and $\gamma>0$. The
following are equivalent.
\begin{enumerate}[label=\textup{(\alph*)},leftmargin=*]
\item $H_s^2\ge\gamma H_s$ on $\mathcal A$ for every $s\in T_K$.
\item $\Phi_\omega(F)\ge0$ for every $\omega$ and every $Z$-even $F\in
\mathcal A^\omega_Q$.
\item $\Phi_\omega(F)\ge0$ for every $\omega$ and every $\operatorname{Stab}
(\omega)$-invariant $F\in\mathcal A^\omega_Q$.
\item \textup{(b)} holds for a dense set of $\omega$, or for $\mu$-almost
every $\omega$.
\end{enumerate}
In particular Theorem~\ref{thm:finite} holds with Definition~\ref{def:patchgap}
restricted to $Z$-even functions, and no condition on a set of boundary
conditions that is not dense can replace~\textup{(b)}.
\end{proposition}

\begin{lemma}\label{lem:fibre}
\begin{enumerate}[label=\textup{(\roman*)},leftmargin=*]
\item For a continuous function $F$ of $\sigma$, $\omega\mapsto\Phi_\omega(F)$
is continuous and depends only on the exterior links of the plaquettes
meeting $V(Q)$.
\item For $g\in\mathcal G_{\rm ext}$, $\Phi_\omega(F\circ g)=\Phi_{g\omega}(F)$.
\item $|\Phi_\omega(F+R)-\Phi_\omega(F)|\le C_m\|R\|_\infty(2\|F\|_\infty+\|R
\|_\infty)$ with $C_m=(m-1)^8+\gamma(m-1)^4$.
\item The exterior graph, whose vertices are the non-interior vertices and
whose edges are the links outside $\Pi_s$, is connected. For a spanning tree
$\mathcal T$ with root $\mathfrak r$ there is a unique continuous $k(\omega)
\in\mathcal G_{\rm ext}$ with $k_{\mathfrak r}=1$ and $(k(\omega)\omega)_e=1$
for $e\in\mathcal T$, and $k(h\omega)=c_{h_{\mathfrak r}}k(\omega)h^{-1}$ for
$h\in\mathcal G_{\rm ext}$.
\item In a compact Lie group there are finitely many elements $g_1,\dots,
g_k$ whose centralisers intersect in the centre; for $\SU(2)$ two
non-commuting elements suffice.
\end{enumerate}
\end{lemma}

\begin{proof}
(i) $\mathcal P^\omega_jF(\sigma)=\int F(\sigma_{-B_j},\tau)p^\omega_j(\tau\mid
\sigma_{-B_j})\,d\tau$ with a density $p^\omega_j$ proportional to $\prod_{p
\cap B_j\ne\emptyset}\varphi(U_p)$, jointly continuous and bounded below, so
$\mathcal P^\omega_jF$ is jointly continuous in $(\sigma,\omega)$; expanding
$\Phi_\omega(F)$ into the terms $\nu^\omega_Q((F-\mathcal P_jF)(F-\mathcal P_{j'}
F))$ and $\nu^\omega_Q((F-\mathcal P_jF)^2)$, each is the integral of a
jointly continuous function against the jointly continuous density of
$\nu^\omega_Q$, and only the plaquettes meeting $V(Q)$ enter.
(ii) $\sigma\mapsto g\sigma$ preserves Haar measure, conjugates every
plaquette variable, hence carries $\nu^\omega_Q$ to $\nu^{g\omega}_Q$ since
$\varphi$ is a class function, and acts link by link, so it intertwines
$\mathcal P^\omega_j$ with $\mathcal P^{g\omega}_j$.
(iii) $\|H^\omega\|\le(m-1)^4$, $\|H(F+R)\|^2-\|HF\|^2=\ip{HR}{H(2F+R)}$,
likewise for the second term, and $\|\cdot\|_{L^2(\nu)}\le\|\cdot\|_\infty$.
(iv) The vertices outside the open box of side $md-1$ form a connected
subgraph of the torus since $L\ge2md$, and each vertex of the lower faces
of $Q$ is joined to its neighbour outside by an exterior link. Let $k_x$ be
the inverse of the holonomy of $\omega$ along the tree path from $\mathfrak r$
to $x$; uniqueness is clear, and $c_{h_{\mathfrak r}}k(\omega)h^{-1}$
trivialises $h\omega$ on $\mathcal T$ and is $1$ at $\mathfrak r$.
(v) For closed subgroups $H'\subsetneq H$ either $\dim H'<\dim H$ or $H'_0
=H_0$ and $H'$ has fewer components. Choose a finite set $S$ for which
$(\dim C(S),|\pi_0C(S)|)$ is lexicographically least; then $C(S\cup\{g\})=C
(S)$ for every $g$, so $C(S)\subset\bigcap_gC(g)=Z\subset C(S)$. In $\SU(2)$
two distinct maximal tori meet in $\{\pm1\}$.
\end{proof}

\begin{proof}[Proof of Proposition~\ref{prop:fibre}]
(b)$\Rightarrow$(a) and (c)$\Rightarrow$(a). Continuous elements are dense in
$\mathcal A$, by averaging continuous approximants over the compact gauge
group of $\Lambda_L$. For continuous $F\in\mathcal A$ and every $\omega$ the
fibre $F_\omega=F(\cdot,\omega)$ lies in $\mathcal A^\omega_Q$, since $\mathcal
G_Q$ fixes $\omega$, and is $\operatorname{Stab}(\omega)$-invariant, since
$F(g\sigma,\omega)=F(g\sigma,g\omega)=F(\sigma,\omega)$; in particular it is
$Z$-even. As in the proof of Theorem~\ref{thm:finite}, $(H_sF)_\omega=H^\omega
F_\omega$, so $\|H_sF\|^2-\gamma\ip F{H_sF}=\int\Phi_\omega(F_\omega)\,d\mu(
\omega)\ge0$, and (a) follows by density.

(b)$\Rightarrow$(c) and (b)$\Rightarrow$(d) are trivial, since
$\operatorname{Stab}(\omega)$-invariant functions are $Z$-even.

(d)$\Rightarrow$(b). Suppose $\Phi_{\omega_0}(F_0)<0$ for a $Z$-even $F_0\in
\mathcal A^{\omega_0}_Q$. Averaging continuous approximants in $L^2(\nu^{
\omega_0}_Q)$ over the compact group generated by $\mathcal G_Q$ and $c_Z$,
which preserves $\nu^{\omega_0}_Q$ by Lemma~\ref{lem:boxinv}(i) and
Lemma~\ref{lem:fibre}(ii), is an orthogonal projection, so $F_0$ may be
taken continuous, $\mathcal G_Q$-invariant pointwise and $Z$-even; these
properties do not depend on $\omega$. By Lemma~\ref{lem:fibre}(i), $\Phi_
\omega(F_0)<0$ on a neighbourhood of $\omega_0$, so the set where (b) holds
is closed, and (b) on a dense set gives (b) everywhere; and since the law
of $\omega$ under $\mu$ has a density bounded below, almost every $\omega$
is a dense set.

(a)$\Rightarrow$(b). Suppose $\Phi_{\omega_0}(F_0)<0$ with $F_0$ continuous,
$\mathcal G_Q$-invariant and $Z$-even, as just arranged. \emph{Generic
stabiliser.} Take $g_1,\dots,g_k$ from Lemma~\ref{lem:fibre}(v), exterior
links $e_1,\dots,e_k$ on no plaquette meeting $V(Q)$, and the tree $\mathcal
T$ avoiding them; this is possible since the exterior graph at distance at
least $2$ from $Q$ contains cycles through any $k$ prescribed links, and
removing links that lie on cycles keeps the graph connected. Let $\omega_1$
agree with $\omega_0$ except that $(k(\omega_0)\omega_1)_{e_i}=g_i$; then $k(
\omega_1)=k(\omega_0)$, since $k$ depends on the tree links only, and $\Phi_{
\omega_1}(F_0)=\Phi_{\omega_0}(F_0)$ by Lemma~\ref{lem:fibre}(i). The slice
point $\hat\omega_1=k(\omega_1)\omega_1$ has $\{x:c_x\hat\omega_1=\hat\omega_1\}
=\bigcap_iC(g_i)=Z$. \emph{Transport.} Put $F_1=F_0\circ k(\omega_1)^{-1}$; it
is $\mathcal G_Q$-invariant, since $\mathcal G_{\rm ext}$ commutes with
$\mathcal G_Q$, and $Z$-even, since $c_z$ commutes with $k(\omega_1)$, and by
Lemma~\ref{lem:fibre}(ii), $\Phi_{\hat\omega_1}(F_1)=\Phi_{\omega_1}(F_0)=-2
\delta<0$. \emph{Neighbourhoods.} By Lemma~\ref{lem:fibre}(i) there is an
open $U_1\ni\hat\omega_1$ with $\Phi_{\omega'}(F_1)\le-\delta$ on $U_1$. By
uniform continuity there is an open $W_0\ni1$ in $G$ such that $\varepsilon=
\sup\{|F_1(c_w\sigma)-F_1(\sigma)|:w\in W_0,\ \sigma\}$ satisfies $C_m
\varepsilon(2\|F_1\|_\infty+\varepsilon)\le\delta/2$. There is an open $U\subset
U_1$ containing $\hat\omega_1$ such that $c_xU\cap U\ne\emptyset$ implies
$x\in W_0Z$: otherwise there are $x_n\notin W_0Z$ and $\omega_n\to\hat\omega_1$
with $c_{x_n}\omega_n\to\hat\omega_1$, a subsequence $x_n\to x$ has $c_x\hat
\omega_1=\hat\omega_1$, so $x\in Z\subset W_0Z$, which is open, against $x_n
\notin W_0Z$. Let $\chi\ge0$ be continuous, supported in $U$, with $\chi(\hat
\omega_1)>0$. \emph{The test function.} Define
\[
F(\sigma,\omega)=\int_G\chi\bigl(c_xk(\omega)\omega\bigr)\,F_1\bigl(c_xk(\omega)
\sigma\bigr)\,dx .
\]
It is bounded, measurable and $\mathcal G_Q$-invariant, and for $h\in
\mathcal G_{\rm ext}$, Lemma~\ref{lem:fibre}(iv) gives $k(h\omega)h=c_{h_{
\mathfrak r}}k(\omega)$, so $F(h\sigma,h\omega)$ is the same integral with $x$
replaced by $xh_{\mathfrak r}$, which is $F(\sigma,\omega)$ by the invariance of
Haar measure; hence $F\in\mathcal A$. \emph{Its fibres.} Put $\hat\omega=k(
\omega)\omega$ and $a(\omega)=\int\chi(c_x\hat\omega)\,dx$, continuous with
$a(\omega_1)>0$. If $a(\omega)=0$, then $F_\omega=0$. Otherwise choose $x_0$
with $c_{x_0}\hat\omega\in U$; whenever $\chi(c_x\hat\omega)\ne0$, both $c_{x_0}
\hat\omega$ and $c_{xx_0^{-1}}(c_{x_0}\hat\omega)$ lie in $U$, so $x=wzx_0$
with $w\in W_0$, $z\in Z$, and by the centrality of $z$ and the $Z$-evenness
of $F_1$, $F_1(c_x\tau)=F_1(c_wc_{x_0}\tau)$ differs from $F_1(c_{x_0}\tau)$ by
at most $\varepsilon$. Hence $F_\omega=a(\omega)\bigl[F_1\circ(c_{x_0}k(\omega))
+R_\omega\bigr]$ with $\|R_\omega\|_\infty\le\varepsilon$, and by
Lemma~\ref{lem:fibre}(ii),(iii) and the homogeneity of $\Phi_\omega$,
\[
\Phi_\omega(F_\omega)\le a(\omega)^2\bigl(\Phi_{c_{x_0}\hat\omega}(F_1)+\delta/2
\bigr)\le-a(\omega)^2\delta/2 .
\]
The set $\{a>0\}$ is open and nonempty, hence of positive $\mu$-mass, so
$\|H_sF\|^2-\gamma\ip F{H_sF}=\int\Phi_\omega(F_\omega)\,d\mu(\omega)<0$,
against (a).

The last claims: the proof of Theorem~\ref{thm:finite} uses only
(b)$\Rightarrow$(a); and (a) forces (b) at every $\omega$.
\end{proof}

At a central boundary condition, one on which $c_x$ acts trivially for all
$x$, $\operatorname{Stab}(\omega)\supset\{c_x\}$ and (c) asks the inequality
for colour-singlet fibres only. The proposition identifies the hypothesis
of Theorem~\ref{thm:finite} with the local Knabe inequality on the torus,
in the sense that the one holds for one $L\ge2md$ if and only if the other
holds for every boundary condition, hence for every such $L$; the passage
from the global inequality $\BP(n_N,\varpi)$ of (D$_\ell$) to that local
inequality, the step that would make the slab hypothesis a consequence of
the cube condition, is the one implication that is not proved here, and
Proposition~\ref{prop:fieldringsall}(v) shows it only in the class of the
field rings.
